% =====================================================================
%  BLOCOS LaTeX PARA INSERÇÃO NO ARTIGO (Lab04S03)
%  Gerados a partir dos gráficos do dashboard (Sprints 1 e 2).
%
%  COMO USAR
%  1. Exporte cada gráfico do dashboard como imagem (PNG/PDF) e salve na
%     pasta de figuras do seu projeto LaTeX com o nome indicado em
%     \includegraphics{...} de cada bloco abaixo.
%  2. Copie os parágrafos + ambientes figure/table para a seção indicada.
%  3. Os números foram apurados sobre o dataset real (N = 202.014 PRs;
%     1.290 repositórios; G1 = IA autônoma, n = 9; G2 = humano, n = 202.005).
%
%  Requer no preâmbulo (já presentes no seu artigo): graphicx, booktabs,
%  adjustbox, float.
% =====================================================================


% #####################################################################
% ###  SEÇÃO 3 (METODOLOGIA) — CARACTERIZAÇÃO DO DATASET             ###
% #####################################################################

\subsection{Caracterização do Dataset}
\label{sec:caracterizacao}

O dataset de análise reúne 202.014 Pull Requests merged, distribuídos em
1.290 repositórios populares, particionados em dois grupos: o Grupo IA (G1),
com 9 PRs produzidos por agentes autônomos, e o Grupo Humano (G2), com
202.005 PRs. Antes da inferência por questão de pesquisa, caracteriza-se a
composição e a tendência central das variáveis de processo, de modo a
contextualizar a forte assimetria amostral e a natureza enviesada das
distribuições. Como as variáveis contínuas são acentuadamente assimétricas e
de cauda longa, adota-se a \textbf{mediana} como medida de tendência central
de referência, reportando-se a média e o intervalo interquartil (IQR) como
apoio.

% --------------------------------------------------------------------
%  Figura: distribuição por linguagem
% --------------------------------------------------------------------
A Figura~\ref{fig:dash-linguagem} apresenta a distribuição dos PRs pelas três
linguagens consideradas. O dataset é dominado por TypeScript, com 98.520 PRs
(48,8\,\%), seguido de Python, com 65.868 PRs (32,6\,\%), e JavaScript, com
37.626 PRs (18,6\,\%). Essa predominância de TypeScript em G2 é relevante
porque contrasta com a composição de G1 --- majoritariamente Python (8 de 9
PRs) ---, o que explica a única diferença estatisticamente robusta observada
no estudo (distribuição de \texttt{language}; ver Seção~\ref{sec:results}) e
reforça que tal diferença é um artefato de composição amostral, não um
indicador de manutenibilidade.

\begin{figure}[H]
\centering
\includegraphics[width=.7\linewidth]{dashboard_dist_linguagem.png}
\caption{Distribuição dos PRs por linguagem de programação no dataset
completo (N = 202.014). TypeScript concentra 48,8\,\% das contribuições.}
\label{fig:dash-linguagem}
\end{figure}

% --------------------------------------------------------------------
%  Figura: distribuição por agente (fingerprint)
% --------------------------------------------------------------------
A Figura~\ref{fig:dash-agente} detalha a procedência das contribuições segundo
o \textit{fingerprint} comportamental do agente de IA associado ao PR. Devin
responde pela maior parcela, com 134.718 PRs (66,7\,\%), seguido de Claude
Code, com 24.769 PRs (12,3\,\%), Cursor, com 3.696 PRs (1,8\,\%), e OpenAI
Codex, com 7 PRs ($\approx$\,0,0\,\%); 38.824 PRs (19,2\,\%) não receberam
atribuição de \textit{fingerprint}. Vale notar que essa identificação opera
sobre todo o dataset --- inclusive sobre PRs de G2 ---, pois o
\textit{fingerprint} sinaliza o ferramental de IA empregado em contribuições
assistidas e confirmadas por humanos, e não apenas as contribuições
plenamente autônomas de G1.

\begin{figure}[H]
\centering
\includegraphics[width=.8\linewidth]{dashboard_dist_agente.png}
\caption{Distribuição dos PRs por agente de IA identificado via
\textit{fingerprint} comportamental \cite{ghaleb:26}. Devin concentra
66,7\,\% das contribuições atribuídas.}
\label{fig:dash-agente}
\end{figure}

% --------------------------------------------------------------------
%  Tabela: tendência central das métricas de processo
% --------------------------------------------------------------------
A Tabela~\ref{tab:dash-central} sumariza a tendência central das principais
variáveis de processo no dataset completo. O contraste sistemático entre
mediana e média evidencia a assimetria positiva das distribuições: em
\texttt{diff\_lines}, por exemplo, a mediana de 71 linhas é muito inferior à
média de 119, indicando uma cauda longa de PRs grandes. O tempo até o merge é
tipicamente curto (mediana de 0,62 dia), porém com média de 4,34 dias puxada
por \textit{outliers}. A baixa proporção de arquivos de teste (mediana 0;
média 0,18) e a presença de revisão formal em 73,5\,\% dos PRs caracterizam o
processo de contribuição predominante.

\begin{table}[H]
\centering
\caption{Tendência central das métricas de processo (dataset completo,
N = 202.014). IQR = intervalo interquartil [P25; P75].}
\label{tab:dash-central}
\begin{adjustbox}{max width=\textwidth}
\small
\begin{tabular}{lrrr}
\toprule
\textbf{Métrica} & \textbf{Mediana} & \textbf{Média} & \textbf{IQR [P25; P75]} \\
\midrule
Linhas alteradas (\texttt{diff\_lines})      & 71   & 118,99 & [28; 172] \\
Arquivos alterados (\texttt{changed\_files}) & 4    & 6,32   & [2; 7] \\
Tempo até o merge (\texttt{days\_to\_merge}) & 0,62 & 4,34   & [0,05; 2,79] \\
Commits por PR (\texttt{commit\_count})      & 3    & 3,70   & [1; 5] \\
Revisões por PR (\texttt{review\_count})     & 1    & 2,32   & [0; 3] \\
Prop. arq. de teste (\texttt{test\_file\_ratio}) & 0 & 0,18 & [0; 0,33] \\
\bottomrule
\end{tabular}
\end{adjustbox}
\end{table}

% --------------------------------------------------------------------
%  Figura: histograma de diff_lines
% --------------------------------------------------------------------
A Figura~\ref{fig:dash-diffhist} mostra o histograma de \texttt{diff\_lines}.
A massa da distribuição concentra-se em PRs pequenos a médios --- a faixa de
0--100 linhas reúne a maioria das contribuições ---, com decaimento acentuado
e cauda longa à direita. Esse formato justifica metodologicamente o uso de
testes não paramétricos (Mann-Whitney U) e da mediana como medida-síntese,
em vez da média, sensível aos \textit{outliers}.

\begin{figure}[H]
\centering
\includegraphics[width=.8\linewidth]{dashboard_hist_diff.png}
\caption{Distribuição (histograma) do número de linhas alteradas por PR
(\texttt{diff\_lines}). Eixo Y: nº de PRs; eixo X: faixas de linhas no diff.
Forte assimetria positiva com cauda longa.}
\label{fig:dash-diffhist}
\end{figure}

% --------------------------------------------------------------------
%  Figura: histograma de days_to_merge
% --------------------------------------------------------------------
A Figura~\ref{fig:dash-dias} apresenta o histograma de \texttt{days\_to\_merge}.
A grande maioria dos PRs é integrada em menos de um dia, refletindo um fluxo
de revisão e merge ágil; ainda assim, observa-se uma cauda de PRs que
permanecem abertos por semanas, responsável pela discrepância entre a mediana
(0,62 dia) e a média (4,34 dias) reportadas na Tabela~\ref{tab:dash-central}.

\begin{figure}[H]
\centering
\includegraphics[width=.8\linewidth]{dashboard_hist_dias.png}
\caption{Distribuição (histograma) do tempo até o merge
(\texttt{days\_to\_merge}). Eixo Y: nº de PRs; eixo X: faixas de dias.
Predomínio de merges em menos de um dia.}
\label{fig:dash-dias}
\end{figure}

% --------------------------------------------------------------------
%  Figura: mediana de diff por linguagem (controle por subgrupo)
% --------------------------------------------------------------------
Por fim, a Figura~\ref{fig:dash-difflang} compara a mediana e a média de
\texttt{diff\_lines} entre as três linguagens. As medianas são notavelmente
homogêneas --- 70 linhas em TypeScript, 73 em Python e 69 em JavaScript ---,
indicando que o tamanho típico das contribuições não varia de forma relevante
entre as linguagens e que a linguagem, isoladamente, não constitui um fator de
confusão para o tamanho dos PRs.

\begin{figure}[H]
\centering
\includegraphics[width=.8\linewidth]{dashboard_diff_por_linguagem.png}
\caption{Mediana e média de \texttt{diff\_lines} por linguagem. As medianas
($\approx$ 70 linhas) são homogêneas entre os subgrupos.}
\label{fig:dash-difflang}
\end{figure}


% #####################################################################
% ###  SEÇÃO 4 (RESULTADOS) — RQ1 e RQ2                              ###
% #####################################################################

% --------------------------------------------------------------------
%  RQ01 — Complexidade Ciclomática (CCM)
% --------------------------------------------------------------------
A Figura~\ref{fig:dash-rq01} sintetiza a resposta à RQ01, exibindo, por meio
de um \textit{box plot} com os pontos individuais sobrepostos, a distribuição
da complexidade ciclomática média (\texttt{ccm\_mean}) por grupo (n = 3 em G1;
n = 32 em G2). A mediana de G1 ($\approx$ 5,38) é substancialmente mais alta
do que a de G2 ($\approx$ 1,25), com tamanho de efeito grande
(Cliff's $\delta = 0{,}71$). Entretanto, o teste de Mann-Whitney U retornou
$p$ bruto $= 0{,}043$ (limítrofe) e, sobretudo, a diferença \textbf{não
sobrevive à correção para múltiplos testes} ($p_{\text{BH}} = 0{,}142$;
$p_{\text{Bonf}} = 0{,}681$). Considerando o tamanho amostral mínimo de G1
(n = 3), o resultado deve ser lido como indício exploratório, e não como
evidência robusta de maior complexidade nas contribuições autônomas.
\textbf{H0 não é rejeitada} ($\alpha = 0{,}05$).

\begin{figure}[H]
\centering
\includegraphics[width=.75\linewidth]{dashboard_rq01_ccm.png}
\caption{RQ01 --- Distribuição da complexidade ciclomática média
(\texttt{ccm\_mean}) por grupo (n = 3 G1; n = 32 G2). A linha central de cada
caixa indica a mediana; pontos representam PRs individuais. G1 exibe mediana
mais alta ($\approx$ 5,38) que G2 ($\approx$ 1,25), porém sem significância
após correção.}
\label{fig:dash-rq01}
\end{figure}

% --------------------------------------------------------------------
%  RQ02 — Densidade de Code Smells / KLoC
% --------------------------------------------------------------------
A Figura~\ref{fig:dash-rq02} responde à RQ02, apresentando a distribuição da
densidade de \textit{code smells} de design por KLoC (\texttt{design\_smell\_%
density}) entre os grupos (n = 9 em G1; n = 50 em G2). Em ambos os grupos a
\textbf{mediana é igual a zero}: a maioria dos PRs não acusa \textit{smells} de
design. As poucas observações positivas --- uma em G1 ($\approx$ 5,52) e três
em G2 ($\approx$ 3,61, 6,10 e 8,03) --- são \textit{outliers} que não alteram a
tendência central. O teste de Mann-Whitney U não indicou diferença
estatisticamente significativa ($p = 0{,}612$; Cliff's $\delta = 0{,}049$,
negligível). \textbf{H0 não é rejeitada} ($\alpha = 0{,}05$). Ressalva
metodológica: nesta execução, os detectores ESLint e Pylint não foram
acionados; a densidade é uma \textit{proxy} obtida via Lizard (funções longas,
NLOC $>$ 50, normalizadas por KLoC), o que não captura violações de convenção
e de acoplamento.

\begin{figure}[H]
\centering
\includegraphics[width=.75\linewidth]{dashboard_rq02_smells.png}
\caption{RQ02 --- Distribuição da densidade de \textit{code smells} de design
por KLoC (\texttt{design\_smell\_density}) por grupo (n = 9 G1; n = 50 G2).
Mediana igual a zero em ambos os grupos; diferença não significativa
($p = 0{,}612$; $\delta = 0{,}049$).}
\label{fig:dash-rq02}
\end{figure}
